% This is part of Mes notes de mathématique
% Copyright (c) 2011-2020
%   Laurent Claessens
% See the file fdl-1.3.txt for copying conditions.

Attention aux conventions. Dans le Frido, un corps peut être réduit à \( \{ 0 \}\) et un idéal premier ne peut pas être \( \{ 0 \}\). Ces conventions ont une série de conséquences un peu partout, par exemple dans la proposition \ref{PROPooRUQKooIfbnQX} où nous parlons d'idéal maximum propre. Comparez par exemple avec \cite{ooWEUDooQybvIx}. Soyez attentifs.

En cas de doutes, nous suivons les conventions de Wikipédia.

%+++++++++++++++++++++++++++++++++++++++++++++++++++++++++++++++++++++++++++++++++++++++++++++++++++++++++++++++++++++++++++
\section{Idéal dans un anneau}
%+++++++++++++++++++++++++++++++++++++++++++++++++++++++++++++++++++++++++++++++++++++++++++++++++++++++++++++++++++++++++++

La définition d'un idéal dans un anneau est la définition~\ref{DefooQULAooREUIU}.

%---------------------------------------------------------------------------------------------------------------------------
\subsection{Idéal maximal}
%---------------------------------------------------------------------------------------------------------------------------

\begin{definition}\label{DEFIdealMax}
Un idéal \( I\) dans un anneau \( A \) est dit \defe{idéal maximal}{idéal!maximal}\index{idéal!maximal} ou idéal maximum si tout idéal \( J \) vérifiant \( I \subset J \subset A \) est soit \( I \), soit \( A \).
\end{definition}

\begin{proposition}[Thème~\ref{THEMEooZYKFooQXhiPD}]     \label{PROPooSHHWooCyZPPw}
    Un idéal \( I\) dans un anneau \( A \) est maximum si et seulement si \( A/I\) est un corps.
\end{proposition}

\begin{proof}
    Soit un idéal maximum \( I\subset A\). Alors les idéaux contenant \( I\) sont \( A\) et \( I\). L'application \( \phi\) de la proposition~\ref{PropIJJIdsousphi} est une bijection, donc l'ensemble des idéaux de \( A/I\) ne contient que deux éléments. Les seuls idéaux de \( A/I\) sont donc \( \{ 0 \}\) et \( A/I\); donc \( A/I\) est un corps par la proposition~\ref{PROPooUOCVooZGAVVk}.

    Dans l'autre sens, c'est la même chose : si \( A/I\) est un corps, il possède exactement deux idéaux (l'idéal \( \{ 0 \}\) et l'anneau entier), donc \( A\) ne contient que deux idéaux contenant \( I \). Donc \( I\) est un idéal maximum.
\end{proof}

%---------------------------------------------------------------------------------------------------------------------------
\subsection{Quotients d'anneaux de polynômes par un idéal}
%---------------------------------------------------------------------------------------------------------------------------
  On a commencé gentiment à parler de polynômes en section~\ref{SEC_minusPolynomes}. Leurs anneaux, et surtout les quotients, fournissent des exemples non triviaux des notions que nous allons aborder dans ce chapitre. Il y en aura aussi dans le chapitre~\ref{CpapMatricesPolyn}.

\begin{definition}[Extension d'anneau par quotient]\label{DEF_ExtensionAnneau}
  Soit \( A \) un anneau, et \( P \in A[X] \) \footnote{Pour rappel de \(A[X] \), c'est en définition~\ref{DefAnneauPolyAX}.}. Si l'on peut déterminer toutes les racines \( \alpha_1, \dots, \alpha_n \) de \( P \) dans un corps contenant \( A \), alors l'anneau \( A[X] / (P) \) est aussi noté \(A[\alpha_1,\dots \alpha_n]\).
\end{definition}

\begin{example}  \label{EX_ConstructionEntiersGauss}
  On considère, dans l'anneau \( \eZ[X] \), le polynôme \( P = X^2 + 1 \). Il est connu que dans le corps des complexes \( \eC \) contenant \( \eZ \), il y a deux racines notées \( i \) et \(-i \), opposées l'une de l'autre. Ainsi, l'anneau quotient \( \eZ[X]/(X^2 + 1) \) peut se récrire \( \eZ[i, -i] \), même encore plus simplement \( \eZ[i] \). Cet anneau, sous-anneau de \( \eC \), est appelé \defe{anneau des entiers de Gauss}{entiers de Gauss}.
\end{example}

\begin{example}  \label{EX_ConstructionZisqrtn}
  Dans le même esprit, toujours dans l'anneau \( \eZ[X] \), on prend le polynôme \( P = X^2 + n \) (avec \( n \in \eN \), \( n \geq 2\). Dans \( \eC \), les racines sont \( i\sqrt n \) et \(-i \sqrt n\), opposées l'une de l'autre. Ainsi, l'anneau quotient \( \eZ[X]/(X^2 + n) \) est l'anneau \( \eZ[i \sqrt n] \).
\end{example}

%+++++++++++++++++++++++++++++++++++++++++++++++++++++++++++++++++++++++++++++++++++++++++++++++++++++++++++++++++++++++++++
\section{Caractéristique}
%+++++++++++++++++++++++++++++++++++++++++++++++++++++++++++++++++++++++++++++++++++++++++++++++++++++++++++++++++++++++++++

\begin{lemmaDef}        \label{LEMDEFooVEWZooUrPaDw}
    Soit l'application
    \begin{equation}
        \begin{aligned}
            \mu\colon \eZ&\to A \\
            n&\mapsto n\cdot 1_A .
        \end{aligned}
    \end{equation}
    \begin{enumerate}
        \item
            C'est un morphisme d'anneaux.
        \item
            Le noyau est un sous-groupe de \( \eZ\)
        \item
            Il existe un unique \( p\in \eZ\) tel que \( \ker(\mu)=p\eZ\).
    \end{enumerate}
    Ce \( p\) est la \defe{caractéristique}{caractéristique!d'un anneau} de \( A\).
\end{lemmaDef}

Par exemple la caractéristique que \( \eQ\) est zéro parce qu'aucun multiple de l'unité n'est nul.

À propos de diagonalisation en caractéristique \( 2\), voir l'exemple~\ref{ExewINgYo}.

\begin{lemma}
    Si \( A\) est de caractéristique nulle, alors \( A\) est infini.
\end{lemma}

\begin{proof}
    En effet, \( \ker\mu=\{0\} \) implique que \( n1_A \neq  m1_A\) dès que \(n \neq m \) et par conséquent \( A\) contient \(\eZ 1_A \), et  est infini.
\end{proof}

\begin{lemma}       \label{LemHmDaYH}
    Si \( p\) est la caractéristique de l'anneau \( A\), alors nous avons l'isomorphisme d'anneaux
    \begin{equation}
         \eZ 1_A\simeq\eZ/p\eZ.
    \end{equation}
\end{lemma}

\begin{proof}
    L'isomorphisme est donné par l'application \( n1_A\mapsto \phi(n)\) si \( \phi\) est la projection canonique \( \eZ\to \eZ/p\eZ\).
\end{proof}

\begin{proposition}     \label{PropGExaUK}
    La caractéristique d'un anneau fini divise son cardinal.
\end{proposition}

\begin{proof}
    Si \( A\) est un anneau, le groupe \( \eZ\) agit sur \( A\) par
    \begin{equation}
        n\cdot a=a+n1_A.
    \end{equation}
    Chaque orbite de cette action est de la forme
    \begin{equation}
        \mO_a=\{ a+n1_A\tq n=0,\ldots, p-1 \}
    \end{equation}
    où \( p\) est la caractéristique de \( A\). Les orbites ont \( p\) éléments et forment une partition de \( A\), donc le cardinal de \( A\) est un multiple de \( p\).
\end{proof}

\begin{lemma}[\cite{ooIBWOooSjOvXd}]        \label{LEMooJQIKooQgukqn}
    Un anneau totalement ordonné est de caractéristique nulle.
\end{lemma}

\begin{proof}
    Le morphisme \( \mu\colon \eZ\to A\), \( n\mapsto n 1_A\) est strictement croissant, en particulier \( \mu(x)\neq \mu(y)\) dès que \( x\neq y\). Donc \( \ker(\mu)=\{ 0 \}\).
\end{proof}

L'ensemble typique de caractéristique \( p\) est \( \eF_p=\eZ/p\eZ\).



\begin{proposition}     \label{Propqrrdem}
    Soit \( A\) un anneau commutatif de caractéristique première \( p\). Alors \( \sigma(x)=x^p\) est un automorphisme de l'anneau \( A\). Nous avons la formule
    \begin{equation}
        (a+b)^p=a^p+b^p
    \end{equation}
    pour tout \( a,b\in A\).
\end{proposition}

\begin{proof}
    Nous utilisons la formule du binôme de la proposition~\ref{PropBinomFExOiL} et le fait que les coefficients binomiaux non extrêmes sont divisibles par \( p\) et donc nuls.
\end{proof}

\begin{proposition} \label{PropFrobHAMkTY}
    Soit \( A\) un anneau commutatif unitaire de caractéristique \( p\). L'application
    \begin{equation}
        \begin{aligned}
            \Frob_A\colon A&\to A \\
            x&\mapsto x^p
        \end{aligned}
    \end{equation}
    est un automorphisme d'anneau unitaire.
\end{proposition}
Nous le nommons le \defe{morphisme de Frobenius}{morphisme!Frobenius}\index{Frobenius!morphisme}. Nous utiliserons aussi les itérés du morphisme de Frobenius : \( \Frob^k\colon x\mapsto x^{p^k}\).

\begin{example}
    Soit à factoriser \( X^p-1\) dans \( \eF_p\). Grâce au morphisme de Frobenius, nous avons immédiatement
    \begin{equation}
        X^p-1=(X-1)^p.
    \end{equation}
\end{example}

%+++++++++++++++++++++++++++++++++++++++++++++++++++++++++++++++++++++++++++++++++++++++++++++++++++++++++++++++++++++++++++
\section{Divisibilité et intégrité, le retour}
%+++++++++++++++++++++++++++++++++++++++++++++++++++++++++++++++++++++++++++++++++++++++++++++++++++++++++++++++++++++++++++
\label{SECAnneauxIntegres}

La définition d'un anneau intègre est la définition~\ref{DEFooTAOPooWDPYmd}.

Nous rappelons (théorème~\ref{ThoBUEDrJ}) que l'anneau \( A\) est intègre si et seulement si \( A[X]\) est intègre.

%---------------------------------------------------------------------------------------------------------------------------
\subsection{Caractéristique d'un anneau intègre}
%---------------------------------------------------------------------------------------------------------------------------

\begin{lemma}       \label{LemCaractIntergernbrcartpre}
    La caractéristique\footnote{Définition~\ref{LEMDEFooVEWZooUrPaDw}.} d'un anneau intègre est zéro ou un nombre premier.
\end{lemma}

\begin{proof}
    Si \( A\) est intègre, alors \( \eZ 1_A\) est a fortiori intègre. Notons \( p \) la caractéristique de \( A \). Si \( p = 0 \), la preuve est finie; supposons donc que \( p \neq 0 \). Alors, l'anneau \( \eZ/p\eZ\) est isomorphe à \( \eZ 1_A\), et est donc intègre. Or, la proposition~\ref{CorZnInternprem} dit que \( \eZ/p\eZ\) est intègre si et seulement si \( p\) est premier, ce qui conclut la preuve.
\end{proof}

\begin{example}
    Il existe des corps dont la caractéristique n'est pas égale au cardinal (contrairement à ce que laisserait penser l'exemple des \( \eZ/p\eZ\)). En effet les matrices \( n\times n\) inversibles sur \( \eF_{3}\) forment un corps qui n'est pas de cardinal trois alors que la caractéristique est \( 3\) :
    \begin{equation}
        \begin{pmatrix}
            1    &       \\
                &   1
            \end{pmatrix}+\begin{pmatrix}
                1    &       \\
                    &   1
                \end{pmatrix}+\begin{pmatrix}
                    1    &       \\
                        &   1
                \end{pmatrix}=0.
    \end{equation}
\end{example}

\begin{example}
    Si \( \eK\) est un corps de caractéristique \( 2\), alors l'égalité \( x=-x\) n'implique pas \( x=0\), vu que \( 2x=0\) est vérifiée pour tout \( x\). Cela se répercute sur un certain nombre de résultats. Par exemple, en caractéristique deux, une forme antisymétrique n'est pas toujours alternée: voir le lemme~\ref{LemHiHNey}.
\end{example}

%---------------------------------------------------------------------------------------------------------------------------
\subsection{Divisibilité et classes d'association}
%---------------------------------------------------------------------------------------------------------------------------
\label{DivisibiliteAnneauxIntegres}

\begin{lemma}\label{LemRmVTRq}
    Si \( A\) est un anneau intègre et si \( a,b\in A\) sont tels que \( a\divides b\) et \( b\divides a\), alors il existe un inversible \( u\in A\) tel que \( a=ub\).
\end{lemma}

\begin{proof}
    Les hypothèses à propos de la divisibilité nous indiquent que \( a=ub\) et \( b=va\) pour certains \( u,v\in A\). Du coup,
    \begin{equation}
        b(1-vu)=0.
    \end{equation}
    Étant donné que \( A\) est intègre, cela montre que \( b=0\) ou \( 1-vu=0\). Si \( b=0\) nous avons immédiatement \( a=0\) et le lemme est prouvé en prenant \( u = 1 \). Si au contraire \( vu=1\), c'est que \( v\) et \( u\) sont inversibles et inverses l'un de l'autre.
\end{proof}

\begin{definition}\label{DefrXUixs}
    On dit de deux éléments \( a,b\in A\) qu'ils sont \defe{associés}{associés!éléments d'un anneau} si ils vérifient les hypothèses du lemme~\ref{LemRmVTRq}; en d'autres termes, $a$ et $b$ sont associés s'il existe un inversible \( u\in A\) tel que \( a=ub\).

    La \defe{classe d'association}{classe d'association}\index{classe!d'association} d'un élément \( a \in A \) est l'ensemble des éléments qui lui sont associés; en d'autres termes, c'est \( a  U(A) \).
\end{definition}

\begin{example}
    Dans \( \eZ[i]\)\footnote{Si vous ne savez pas ce que c'est, c'est que vous n'avez pas lu le chapitre depuis le début: c'est bien, mais quand même, allez voir l'exemple~\ref{EX_ConstructionEntiersGauss}\dots}, les inversibles sont \( \pm 1\) et \( \pm i\). Donc les éléments associés à \( z\) sont \( z\), \( -z\), \( iz\) et \( -iz\).

    Notons au passage que la notion de divisibilité dans \( \eZ[i]\) n'est pas immédiatement intuitive. En effet bien que \( 7\) ne soit pas divisible par \( 2\) (ni dans \( \eZ\) ni dans \( \eZ[i]\)), le nombre \( 7+6i\) est divisible par \( 2+i\) dans \( \eZ[i]\). En effet :
    \begin{equation}
        (2+i)(4+i)=7+6i.
    \end{equation}
\end{example}

\begin{probleme}
    Est-ce que quelqu'un connaît un anneau contenant \( \eZ\) dans lequel \( 7\) est divisible en \( 2\) ?

    Peut-être \( \eZ\) étendu par tous les \( 1/2^n\) ? En tout cas, on a besoin de "7/2".
\end{probleme}

%---------------------------------------------------------------------------------------------------------------------------
\subsection{PGCD et PPCM}
%---------------------------------------------------------------------------------------------------------------------------

Pour le théorème de Bézout et autres opérations avec des modulo, voir le thème~\ref{THEMEooNRZHooYuuHyt}. Le pgcd et le ppcm sont définis en \ref{DefrYwbct}.

\begin{lemma}
    Soient \( A\) un anneau intègre et \( S\subset A\). Si \( \delta\) est un PGCD de \( S\), alors l'ensemble des PGCD de \( S\) est la classe d'association de \( \delta\).

    De la même façon si \( \mu\) est un PPCM de \( S\), alors l'ensemble des PPCM de \( S\) est la classe d'association de \( \mu\).
\end{lemma}

\begin{proof}
    Soient \( \delta\) un PGCD de \( S\) et \( u\) un inversible dans \( A\). Si \( x\in S\) nous avons \( \delta\divides x\) et donc \( x=a\delta\). Par conséquent \( x=au^{-1}u\delta\) et donc \( u\delta\) divise \( x\). De la même manière, si \( d\) divise \( x\) pour tout \( x\in S\), alors \( d\) divise \( \delta\) et donc \( \delta=ad\) et \( u\delta=aud\), ce qui signifie que \( d\) divise \( u\delta\).

    Dans l'autre sens nous devons prouver que si \( \delta'\) est un autre PGCD de \( S\), alors il existe un inversible \( u\in A\) tel que \( \delta'=u\delta\). Vu que \( \delta'\) divise \( x\) pour tout \( x\in S\), nous avons \( \delta'\divides \delta\), et symétriquement nous trouvons \( \delta\divides\delta'\). Par conséquent (lemme~\ref{LemRmVTRq}), il existe un inversible \( u\) tel que \( \delta=u\delta'\).

    Le même type de raisonnement tient pour le PPCM.
\end{proof}

Si \( \delta\) est un PGCD de \( S\), nous dirons \emph{par abus de langage} que \( \delta\) est \emph{le} PGCD de \( S\), gardant en tête qu'en réalité toute sa classe d'association est PGCD. Nous noterons aussi, toujours par abus que \( \delta=\pgcd(S)\).

\begin{remark}
    La classe d'association d'un élément n'est pas toujours très grande. Les inversibles dans \( \eZ\) étant seulement \( \pm 1\), nous pouvons obtenir l'unicité du PGCD et du PPCM en imposant qu'ils soient positifs: c'est ce que nous avons observé avec le lemme \ref{LEMooBJVJooFyuFeN}.

    Pour les polynômes, nous obtenons l'unicité en demandant que le PGCD soit unitaire\footnote{Définition~\ref{DEF_polynUnitaire}.}.

    Dans les cas pratiques, il y a donc en réalité peu d'ambiguïté à parler du PGCD ou du PPCM d'un ensemble.
\end{remark}

%--------------------------------------------------------------------------------------------------------------------------- 
\subsection{Élément irréductible}
%---------------------------------------------------------------------------------------------------------------------------

\begin{definition}[Élément irréductible\cite{ooWUNIooXKxRya}]  \label{DeirredBDhQfA}
    Un élément d'un anneau commutatif est \defe{irréductible}{irréductible!dans un anneau} si il n'est ni inversible, ni le produit de deux éléments non inversibles.
\end{definition}

\begin{normaltext}
    Nous allons voir dans la section \ref{SECooSWGKooEeOZTO} que le concept d'élément irréductible n'est vraiment utile que dans le cas des anneaux intègres.
\end{normaltext}

\begin{example}
    Un corps n'a pas d'éléments irréductibles parce qu'à part zéro tous les éléments sont inversibles. Mais \( 0\) n'est pas irréductible parce qu'il peut être écrit comme produit d'éléments non inversibles : \( 0=0\cdot 0\).
\end{example}

\begin{example}
    Les éléments irréductibles de l'anneau \( \eZ\) sont les nombres premiers. En effet les seuls inversibles de \( \eZ\) sont \( \pm 1\). Si \( p\) est premier et \( p=ab\) avec \( a,b\in \eZ\), alors nous avons soit \( a=\pm 1\) soit \( b=\pm 1\).
\end{example}

%+++++++++++++++++++++++++++++++++++++++++++++++++++++++++++++++++++++++++++++++++++++++++++++++++++++++++++++++++++++++++++
\section{Des anneaux particuliers}
%+++++++++++++++++++++++++++++++++++++++++++++++++++++++++++++++++++++++++++++++++++++++++++++++++++++++++++++++++++++++++++

%--------------------------------------------------------------------------------------------------------------------------- 
\subsection{Anneau factoriel}
%---------------------------------------------------------------------------------------------------------------------------

\begin{definition}[Anneau factoriel]        \label{DEFooVCATooPJGWKq}
    Un anneau commutatif \( A\) est \defe{factoriel}{factoriel!anneau}\index{anneau!factoriel} s'il vérifie les propriétés suivantes.
    \begin{enumerate}
        \item
            L'anneau \( A\) est intègre (pas de diviseurs de zéro).
        \item
            Si \( a\in A\) est non nul et non inversible, alors il admet une décomposition en facteurs irréductibles: \( a=p_1\ldots p_k\) où les \( p_i\) sont irréductibles.
        \item
            Si \( a=q_1\ldots q_m\) est une autre décomposition de \( a\) en irréductibles, alors \( m=k\) et il existe une permutation\footnote{Définition~\ref{DEFooJNPIooMuzIXd}.} \( \sigma\in S_k\) telle que \( p_i\) et \( q_{\sigma(i)}\) soient associés\footnote{Définition~\ref{DefrXUixs}.}.
    \end{enumerate}
\end{definition}

Un anneau factoriel permet de caractériser le \( \pgcd\) et le \( \ppcm\) de nombres.

\begin{proposition}
Soit une famille \( \{ a_n \}\) d'éléments de \( A\) qui se décomposent en irréductibles comme
\begin{equation}
    a_i=\prod_k p_k^{\alpha_{k,i}}.
\end{equation}
Alors
\begin{equation}
    \pgcd\{ a_n \}=\prod_k p_k^{\min_i\{ \alpha_{k,i} \}}.
\end{equation}

De plus le PGCD est :
\begin{enumerate}
    \item
        Un multiple de tous les diviseurs communs des \( a_i\).
    \item
        Unique pour cette propriété à multiple près par un inversible\quext{Soyez prudent avec cette affirmation : je n'en n'ai pas de démonstrations sous la main et ne suis pas certain que ce soit vrai.}.
\end{enumerate}

\end{proposition}

De la même manière,
\begin{equation}
    \ppcm\{ a_n \}=\prod_kp_k^{\max_i\{ \alpha_{k,i} \}}.
\end{equation}
Un anneau factoriel a une relation de préordre partiel\index{ordre!sur un anneau factoriel} donnée par \( a<b\) si \( a\) divise \( b\). En termes d'idéaux, cela donne l'ordre inverse de celui de l'inclusion\footnote{Voir proposition~\ref{PropDiviseurIdeaux}.} : \( a<b\) si et seulement si \( (b)\subset (a)\).

\begin{example} \label{EXooCWJUooCDJqkr}
    L'anneau \( \eZ[i\sqrt{3}]\) n'est pas factoriel parce que \( 4\) a au moins deux décompositions distinctes en irréductibles :
    \begin{equation}
        4=2\cdot 2,
    \end{equation}
    et
    \begin{equation}
        4=(1+i\sqrt{3})(1-i\sqrt{3}).
    \end{equation}
\end{example}

Nous allons voir dans l'exemple~\ref{ExeDufyZI} que \( \eZ[i\sqrt{2}]\) est factoriel parce qu'il sera euclidien.

%--------------------------------------------------------------------------------------------------------------------------- 
\subsection{Anneau principal et idéal premier}
%---------------------------------------------------------------------------------------------------------------------------
%///////////////////////////////////////////////////////////////////////////////////////////////////////////////////////////
\subsubsection{Idéal premier}
%///////////////////////////////////////////////////////////////////////////////////////////////////////////////////////////

\begin{definition}      \label{DEFooAQSZooVhvQWv}
    Nous disons qu'un idéal \( I\) dans \( A\) est \defe{premier}{premier!idéal} si \( I\) est strictement inclus dans \( A\) et si pour tout \( a,b\in A\) tels que \( ab\in I\) nous avons \( a\in I\) ou \( b\in I\).
\end{definition}

\begin{lemma}       \label{LEMooYRPBooYxXXsi}
    L'idéal nul (réduit à \( \{ 0 \}\)) est premier si et seulement si \( A\) est intègre.
\end{lemma}

\begin{proof}
    En deux sens.
    \begin{subproof}
    \item[Si \( \{ 0 \}\) est premier]
        Alors \( A\neq \{ 0 \}\) parce que \( I=\{ 0 \}\) est propre (définition d'idéal premier).
        
        De plus, si \( ab=0\), alors \( ab\in I\) qui est un idéal premier. Donc soit \( a\) soit \( b\) est dans \( I\), c'est-à-dire que soit \( a\) soit \( b\) est nul. Donc \( A\) est intègre.

    \item[Si \( A\) est intègre]

        Alors \( A\neq \{ 0 \}\) et l'idéal \( I=\{ 0 \}\) est strictement inclus dans \( A\). Si \( ab\in I\), alors \( ab=0\) et comme \( A\) est intègre, soit \( a\) soit \( b\) est nul, c'est-à-dire appartient à \( I\).
    \end{subproof}
\end{proof}

\begin{proposition}[\cite{ooWEUDooQybvIx, ooOYKZooOJBDHS}]      \label{PROPooRUQKooIfbnQX}   \label{PROPooHABIooBZZQMj}
    Soit un anneau commutatif et un idéal \( I\) dans \( A\).
    \begin{enumerate}
        \item       \label{ITEMooUGBTooOGrnWl}
            \( I\) est un idéal premier si et seulement si \( A/I\) est un anneau intègre.
        \item   \label{ITEMooGLXSooUjINqR}
            \( I\) est un idéal maximal si et seulement si \( A/I\) est un corps.
        \item       \label{ITEMooTFFQooOUajFw}
            Tout idéal maximal propre est premier.
    \end{enumerate}
\end{proposition}


\begin{proof}
    Il n'y a que \ref{ITEMooUGBTooOGrnWl} et \ref{ITEMooTFFQooOUajFw} à vraiment prouver, parce que \ref{ITEMooGLXSooUjINqR} est la proposition \ref{PROPooSHHWooCyZPPw}. Nous allons néanmoins en donner une autre preuve, plus constructive.
    \begin{subproof}
        \item[ \ref{ITEMooUGBTooOGrnWl} -- \( I\) premier implique \( A/I\) intègre]
            Évacuons le cas trivial pour être sûr. Si \( I=\{ 0 \}\) alors \( A\) est intègre par le lemme \ref{LEMooYRPBooYxXXsi}. Donc \( A/I=A/\{ 0 \}=A\) est intègre également.

            Soient \( a,b\in A\) tels que \( [a][b]=[0]\). Donc \( [ab]=[0]\), c'est-à-dire \( ab\in I\). Vu que \( I\) est un idéal premier nous avons \( a\in I\) ou \( b\in I\), c'est-à-dire \( [a]=0\) ou \( [b]=0\); nous en déduisons que \( A/I\) est un anneau intègre.

        \item[ \ref{ITEMooUGBTooOGrnWl} -- \( A/I\) intègre implique \( I\) premier]
            Soit \( ab\in I\). Alors \( [ab]=0\), ce qui signifie que \( [a][b]=0\) donc que \( [a]=0\) ou que \( [b]=0\) parce que \( A/I\) est intègre. Mais la condition \( [a]=0\) signifie \( a\in I\), et \( [b]=0\) signifie \( b\in I\). Nous avons donc prouvé que soit \( a\) soit \( b\) est dans \( I\), c'est-à-dire que \( I\) est premier.

        \item[\ref{ITEMooGLXSooUjINqR} -- \( I\) maximum implique \( A/I\) corps]

            Nous devons montrer que tout élément non nul de \( A/I\) est inversible. Un élément non nul de \( A/I\) est \( [x]\) avec \( x\in A\setminus I\). 
            
            Nous considérons \( J=Ax+I\), qui est un idéal parce que pour tout \( a\in A\), \( aAx+aI\in Ax+I\). Mais comme \( I\) est maximal, \( J=I\) ou \( J=A\).

            Si \( J=I\), nous aurions que pour tout \( a\in A\) et pour tout \( i\in I\), \( ax+i\in I\). En particulier pour \( a=1\) et \( i=0\) nous aurions \( x\in I\), ce qui est contraire à l'hypothèse faite sur \( x\).

            Donc \( J=A\). En particulier, \( 1\in J\), c'est-à-dire qu'il existe \( a\in A\) et \( i\in I\) tels que \( ax+i=1\). En passant aux classes, \( [ax]=1\), c'est-à-dire \( [a][x]=1\) qui signifie que \( [a]\) est un inverse de \( [x]\) dans \( A/I\).

        \item[\ref{ITEMooGLXSooUjINqR} -- \( A/I\) est corps implique \( I\) maximum]

            Si \( x\in A\setminus I\), il faut prouver que tout idéal contenant \( I\) et \( x\) est \( A\).

            Un idéal contenant \( I\) et \( x\) doit contenir l'idéal \( J=Ax+I\). Vu que \( x\notin I\), nous avons \( [x]\neq 0\) dans \( A/I\). Donc \( [x] \) est inversible et il existe \( a\in A\) tel que \( [ax]=[A]\). C'est-à-dire que $ax-1\in I$. Nous avons alors
            \begin{equation}
                1=ax+\underbrace{(1-ax)}_{\in I}.
            \end{equation}
            C'est-à-dire que \( 1\in Ax+I\) et donc \( Ax+I=A\).
    \end{subproof}
    Il reste à prouver \ref{ITEMooTFFQooOUajFw}: tout idéal maximal propre est premier. 

    Si \( I\) est maximal, \( A/I\) est un corps par le point \ref{ITEMooGLXSooUjINqR}, et vu que \( I\) est propre, le corps \( A/I\) n'est pas réduit à \( \{ 0 \}\). Donc le lemme \ref{LemAnnCorpsnonInterdivzer} dit que \( A/I\) est un anneau intègre. Le point \ref{ITEMooUGBTooOGrnWl} dit alors que \( I\) est un idéal premier.
\end{proof}

\begin{remark}
    Vu qu'un corps peut être réduit à \( \{0\}\), dans \ref{ITEMooGLXSooUjINqR}, l'idéal peut être \( A\). Mais pas dans \ref{ITEMooTFFQooOUajFw}, parce qu'un idéal premier est propre, ça fait partie de la définition \ref{DEFooAQSZooVhvQWv}.
\end{remark}

%///////////////////////////////////////////////////////////////////////////////////////////////////////////////////////////
\subsubsection{Idéal principal, anneau principal}
%///////////////////////////////////////////////////////////////////////////////////////////////////////////////////////////
\begin{definition}      \label{DEFooMZRKooBPLAWH}
    Un idéal \( I\) dans \( A\) est \defe{principal à gauche}{idéal!principal!à gauche} s'il existe \( a\in I\) tel que \( I= A a\). Il est \defe{principal à droite}{idéal!principal!à droite} s'il existe \( a\in I\) tel que \( I=a A\). Nous disons qu'il est \defe{principal}{principal!idéal} s'il est principal à gauche et à droite.
\end{definition}

\begin{definition}          \label{DEFooGWOZooXzUlhK}
    Un anneau est \defe{principal}{principal!anneau} si
    \begin{enumerate}
        \item
            il est commutatif et intègre
        \item
            tous ses idéaux sont principaux.
    \end{enumerate}
\end{definition}

Souvent pour prouver qu'un anneau est principal, nous prouvons qu'il est euclidien (définition~\ref{DefAXitWRL}) et nous utilisons la proposition~\ref{Propkllxnv} qui dit qu'un anneau euclidien est principal.

Une manière de prouver qu'un anneau n'est pas principal est de prouver qu'il n'est pas factoriel, théorème~\ref{THOooANCAooBChmwp}.

\begin{proposition}[Thème~\ref{THEMEooZYKFooQXhiPD}, \cite{MonCerveau}] \label{PropomqcGe}
    Soit \( A\) un anneau principal\footnote{Définition \ref{DEFooGWOZooXzUlhK}.} qui n'est pas un corps. Pour un idéal propre \( I\) de \( A\), les conditions suivantes sont équivalentes :
    \begin{enumerate}
        \item       \label{ITEMooNOVFooEHtcwE}
            \( I\) est un idéal maximal\footnote{Définition \ref{DEFIdealMax}.};
        \item       \label{ITEMooMQWVooNocVEU}
            \( I\) est un idéal premier non nul\footnote{Définition \ref{DEFooAQSZooVhvQWv}.};
        \item       \label{ITEMooJBXGooEISNuW}
            il existe \( p\) irréductible\footnote{Définition \ref{DeirredBDhQfA}.} dans \( A\) tel que \( I=(p)\).
    \end{enumerate}
\end{proposition}

\begin{proof}
    En plusieurs implications.
    \begin{subproof}
        \item[\ref{ITEMooNOVFooEHtcwE} implique~\ref{ITEMooMQWVooNocVEU}]

            Par hypothèse, \( I\) est un idéal propre, de plus il n'est pas égal à \( \{ 0 \}\), parce que lorsque \( A\) et \( \{ 0 \} \) sont les seuls idéaux, nous avons un corps (proposition~\ref{PROPooUOCVooZGAVVk}). Étant donné que \( I\) est un idéal maximal, le quotient \( A/I\) est un corps par la proposition~\ref{PROPooSHHWooCyZPPw}. Le quotient est donc, comme tout corps, un anneau intègre (exemple~\ref{EXooMXNTooZaRPPi}), et la proposition~\ref{PROPooRUQKooIfbnQX} dit qu'alors, l'idéal \(I \) est premier.


        \item[\ref{ITEMooMQWVooNocVEU} implique~\ref{ITEMooJBXGooEISNuW}]

            Maintenant \( I\) est un idéal premier non réduit à \( \{ 0 \}\). Vu que \( A\) est un anneau principal, il existe \( x\in A\) tel que \( I=(x)\). Nous devons prouver que \( x\) peut être choisi irréductible; et nous allons faire plus : nous allons prouver que \( x\) ne peut être que irréductible\quext{ça me semble un peu trop facile. Lisez attentivement, et écrivez-moi pour dire si vous êtes d'accord ou pas.}.

            Supposons que \( x\) ne soit pas irréductible. Alors il existe \( a,b\in A\) non inversibles tels que \( x=ab\). Si \( a\in (x)\) alors il existe \( k\in A\) tel que \( a=xk\), et en particulier, \( a=abk\), c'est-à-dire \( 1=bk\) (parce que \( A\) est principal et donc intègre). Cela signifie que \( b\) est inversible alors que nous avions dit qu'il ne l'était pas. Nous en déduisons que \( a\) n'est pas dans \( (x)\). On montre de manière similaire que \( b\) n'est pas dans \( (x)\) non plus.

            Nous nous retrouvons donc avec \( a,b\in A\) tel que \( ab\in I\) sans que ni \( a\) ni \( b\) sont soient dans \( I\). Cela contredit le fait que \( I\) soit un idéal premier. En conclusion, \( x\) est irréductible.

        \item[\ref{ITEMooJBXGooEISNuW} implique~\ref{ITEMooNOVFooEHtcwE}]

            Nous avons \( I=(p)\) avec \( p\) irréductible dans \( A\). Supposons que \( J\) est un idéal différent de \( A\) contenant \( I\). Vu que \( A\) est principal, il existe \( y\in A\) tels que \( J=(y)\). En particulier \( p\in J\), donc \( p=ay\) pour un certain \( a\in A\). Mais \( p\) est irréductible, donc soit \( a\) est inversible, soit \( y\) est inversible. Si \( y\) est inversible, alors \( J=A\), ce qui est exclu. Si \( a\) est inversible, alors \( (y)=(p)\), et \( I=J\).
    \end{subproof}
\end{proof}

\begin{normaltext}
    Dans la proposition \ref{PropomqcGe}, l'hypothèse d'idéal propre est importante. En effet dans le cas \( I=A\), nous avons évidemment que \( I\) est un idéal maximum. Mais \( A\) n'est d'abord pas un idéal premier parce qu'un idéal premier doit être strictement inclus dans l'anneau. Et ensuite, \( A\) est en général loin d'être garanti d'être égal à \( (p)\) pour un de ses éléments \( p\).
\end{normaltext}

\begin{proposition}     \label{PropoTMMXCx}
    Soit \( A \) un anneau principal, et soit \( p \in A \) un élément irréductible. Alors
    \begin{enumerate}
        \item
            \( (p)\) est un idéal maximum.
        \item       \label{ITEMooKPJQooWuPZXS}
            \( A/(p)\) est un corps.
    \end{enumerate}
\end{proposition}

\begin{proof}
    Nous notons \( I=(p)\). Soit un idéal \( J\) contenant \( I\). Vu que \( A\) est principal, \( J\) aussi est monogène : \( J=(q)\). Mais comme \( p\) est dans \( I\) qui est dans \( J\), il existe \( a\in A\) tel que \( p=qa\).

    Vu que \( p\) est irréductible, soit \( q\) soit \( a\) est inversible. Si \( q\) est inversible, alors \( J=A\). Si \( a\) est inversible, alors nous avons \( p=qa\), donc \( q=pa^{-1}\), ce qui signifie que \( q\in(p)\) et donc que \( J=I\).

    Cela prouve que \( (p)\) est un idéal maximum.

    Le fait que \( A/(p)\) soit un corps est maintenant la proposition~\ref{PROPooSHHWooCyZPPw}.
\end{proof}

\begin{example}
    L'anneau \( \eZ\) est principal parce qu'il est intègre et que ses seuls idéaux sont les \( n\eZ\) qui sont principaux : \( n\eZ\) est engendré par \( n\).
\end{example}

\begin{example}[Les idéaux de $\eZ/n\eZ$]       \label{EXooCJRPooYkWdyr}

    Les idéaux de \( \eZ/n\eZ\) sont principaux, mais l'anneau \( \eZ/n\eZ\) n'est pas principal lorsque \( n\) n'est pas premier. Nous allons voir ça.

    \begin{subproof}
        \item[Les idéaux de \( \eZ/n\eZ\) sont principaux]

            Soit un idéal \( S\) dans \( \eZ/n\eZ\). Nous considérons la projection canonique \( \phi\colon \eZ\to \eZ/n\eZ\). La proposition~\ref{PropIJJIdsousphi} dit que  \( S=\phi(J)\) où \( J\) est un idéal de \( \eZ\) contenant \( n\eZ\). Mais le corolaire~\ref{CORooLINXooBlUKPG} nous dit qu'alors \( J=m\eZ\) pour un certain \( m\). Pour que \( m\eZ\) contienne \( n\eZ\), il faut que \( m\) divise \( n\).

            Bref, \( S=\phi(m\eZ)\) avec \( m\divides n\). Nous montrons maintenant que \( S\) est engendré par \( [m]_n\). D'abord, l'élément \( [m]_n\) est bien dans \( \phi(m\eZ)\). Ensuite un élément de \( \phi(m\eZ)\) est de la forme
            \begin{equation}
                [km]_n=k[m]_n\in ([m]_n).
            \end{equation}
            Donc \( S\subset ([m]_n)\). Et l'inclusion dans l'autre sens est tout aussi immédiate : un élément de \( ([m]_n)\) est de la forme
            \begin{equation}
                k[m]_n=[km]_n=\phi(km)\in \phi(m\eZ).
            \end{equation}

        \item[Si \( n\) n'est pas premier, \( \eZ/n\eZ\) n'est pas principal]

            Le fait est que lorsque \( n\) n'est pas premier, \( \eZ/n\eZ\) n'est pas intègre (corolaire~\ref{CorZnInternprem}).

        \item[Moralité]

            Un anneau comme \( \eZ/6\eZ\) est un anneau dont tous les idéaux sont principaux, mais qui n'est pas principal.

    \end{subproof}
\end{example}

\begin{example}
    Nous verrons dans la proposition~\ref{PROPooVWRPooGQMenV} que l'anneau des fonctions holomorphes sur un compact de \( \eC\) est principal.
\end{example}

\begin{definition}      \label{DEFooXSPFooPumQSy}
Nous disons que deux éléments d'un anneau principal sont \defe{premiers entre eux}{premier!deux éléments d'un anneau principal} si leur PGCD est \( 1\).
\end{definition}

\begin{theorem}\index{théorème!chinois!anneau principal}        \label{ThofPXwiM}
    Si \( A\) est un anneau principal et si \( p\) et \( q\) sont premiers entre eux dans \( A\), alors on a l'isomorphisme d'anneaux
    \begin{equation}
        A/pqA\simeq A/pA\times A/qA.
    \end{equation}
\end{theorem}
% TODO : trouver une preuve. Je parie que recopier la même que celle dans Z fonctionne très bien.

%///////////////////////////////////////////////////////////////////////////////////////////////////////////////////////////
\subsubsection{Bézout}
%///////////////////////////////////////////////////////////////////////////////////////////////////////////////////////////
\begin{theorem}[\cite{XPXxPl}]
    Toute partie \( S\) d'un anneau principal admet un PGCD et un PPCM. De plus
    \begin{equation}
        \begin{aligned}[]
            \delta=\pgcd(S)&\Leftrightarrow (\delta)=\sum_{s\in S}(s)\\
            \mu=\ppcm(S)&\Leftrightarrow (\mu)=\bigcap_{s\in S}(s)
        \end{aligned}
    \end{equation}
\end{theorem}

\begin{proof}
    Vu que l'anneau \( A\) est principal, tous ses idéaux sont principaux et donc engendrés par un seul élément. En particulier il existe \( \delta,\mu\in A\) tels que
    \begin{subequations}
        \begin{align}
            (\delta)&=\sum_{s\in S}(s)\\
            (\mu)&=\bigcap_{s\in S}(s)
        \end{align}
    \end{subequations}
    \begin{subproof}
    \item[PGCD]
        Montrons ce que \( \delta\) est un PGCD de \( S\). Pour tout \( x\in S\), nous avons \( (x)\subset (\delta)\), et donc \( \delta\divides x\). Par ailleurs si \( d\divides x\) pour tout \( x\in S\), nous avons \( (x)\subset (d)\) et donc
        \begin{equation}
            \sum_{x\in S}(x)\subset (d),
        \end{equation}
        puis \( (\delta)\subset (d)\) et finalement \( d\divides \delta\).
        \item[PPCM]
            Si \( x\in S\) nous avons \( (\mu)\subset (x)\) et donc \( x\divides \mu\). D'autre part si \( x\divides m\) pour tout \( x\in S\), alors \( (m)\subset (x)\) et donc \( (m)\subset(\mu)\), finalement \( \mu\divides m\).
    \end{subproof}
\end{proof}

\begin{corollary}[Théorème de Bézout\cite{XPXxPl}]\index{Bézout!anneau principal}\label{CorimHyXy}
    Soit un anneau principal \( A\). Deux éléments \( a,b\in A\) sont premiers entre eux si et seulement s'il existe un couple \( (u, v)\in A^2 \) tel que
    \begin{equation}
        ua+vb=1.
    \end{equation}
    À la place de \( 1\) on aurait pu écrire n'importe quel inversible.
\end{corollary}
\index{anneau!principal}

\begin{proof}
    Pour cette preuve, nous allons écrire \( \pgcd(a,b)\) l'ensemble de PGCD de \( a\) et \( b\), c'est-à-dire la classe d'association d'un PGCD.

    Si \( a\) et \( b\) sont premiers entre eux, alors
    \begin{equation}
        1\in\pgcd(a,b)=\sum_{x=a,b}(x)=(a)+(b).
    \end{equation}

    À l'inverse, si nous avons \( ua+vb=1\), alors \( 1\in (a)+(b)\), mais vu que \( (a)+(b)\) est un idéal principal, \( (1)=(a)+(b)\) et donc \( 1\in \pgcd(a,b)\).
\end{proof}

Le lemme de Gauss est une application immédiate de Bézout. Il y aura aussi un lemme de Gauss à propos de polynômes (lemme~\ref{LemEfdkZw}), et une généralisation directe au théorème de Gauss, théorème~\ref{ThoLLgIsig}.
\begin{lemma}[\href{http://ljk.imag.fr/membres/Bernard.Ycart/mel/ar/node6.html}{lemme de Gauss}]    \label{LemSdnZNX}
    Soit \( A\) un anneau principal et \( a,b,c\in A\) tels que \( a\) divise \( bc\). Si \( a\) est premier avec \( c\), alors \( a\) divise \( b\).
\end{lemma}
\index{lemme!Gauss!dans un anneau principal}

\begin{proof}
    Vu que \( a\) est premier avec \( c\), nous avons \( \pgcd(a,c)=1\) et Bézout (\ref{ThoBuNjam}) nous donne donc \( s,t\in \eA\) tels que \( sa+tc=1\). En multipliant par \( b\),
    \begin{equation}
        sab+tbc=b.
    \end{equation}
    Mais les deux termes du membre de gauche sont multiples de \( a\) parce que \( a\) divise \( bc\). Par conséquent \( b\) est somme de deux multiples de \( a\) et donc est multiple de \( a\).
\end{proof}
Un cas usuel d'utilisation est le cas de \( A=\eN^*\).

%///////////////////////////////////////////////////////////////////////////////////////////////////////////////////////////
\subsubsection{Élément premier}
%///////////////////////////////////////////////////////////////////////////////////////////////////////////////////////////

\begin{definition}[\cite{ooWBLYooLYwALS}]       \label{DEFooZCRQooWXRalw}
    Soit un anneau commutatif \( A\). Un élément \( p\in A\) est \defe{premier}{élément premier} si il est
    \begin{enumerate}
        \item
            non nul et non inversible, et
        \item       \label{ITEMooPMTTooCVHPIm}
            tel que, s'il divise un produit \( ab\), alors il divise soit \( a\) soit \( b\) (ou les deux).
    \end{enumerate}
\end{definition}

\begin{proposition}[\cite{ooTGPAooQTbamu}]     \label{PROPooWMNPooZdvOBt}
    Dans un anneau intègre\footnote{Si pas intègre, voir l'exemple \ref{EXooEIUEooCZCPMC}.} tout élément premier est irréductible\footnote{Toutes les définitions dans le thème \ref{THEMEooVIQIooOcFAQS}. Ici, penser à la définition~\ref{DeirredBDhQfA}}.
\end{proposition}
    
\begin{proof}
    Soit \( p\), un élément premier dans un anneau intègre \( A\). Alors  \( p\) n'est pas inversible, cela fait partie de la définition d'un élément premier. Vérifions par l'absurde qu'il n'est pas un produit d'inversibles.
    
    Soient \( a,b\in A\) tels que \( p=ab\). Par le point \ref{ITEMooPMTTooCVHPIm} de la définition \ref{DEFooZCRQooWXRalw}, \( p\) divise soit \( a\) soit \( b\). Supposons que \( p\) divise \( a\). Alors il existe \( x\in A\) tel que \( a=px\). En remettant dans \( p=ab\) nous avons :
            \begin{equation}        \label{EQooPYBGooLFHMJZ}
                p=pxb.
            \end{equation}
            Mais l'anneau est intègre et permet donc des simplifications par tout élément non nul\footnote{Proposition~\ref{DEFooTAOPooWDPYmd}}. La relation \ref{EQooPYBGooLFHMJZ} donne donc 
            \begin{equation}
                1=xb,
            \end{equation}
            ce qui signifie que \( b\) est inversible.

            Un travail similaire montre que \( a\) est inversible si \( p\) divise \( b\).
\end{proof}

\begin{example}
    Si nous avons l'égalité \( 7=ab\) dans \( \eZ\), alors soit \( a\) soit \( b\) vaut \( 1\) et est donc inversible.
\end{example}

Sur un anneau non intègre, la notion d'élément premier n'est pas aussi intéressante que sur un anneau intègre. Par exemple la proposition \ref{PROPooWMNPooZdvOBt} devient fausse.

\begin{example}     \label{EXooEIUEooCZCPMC}
    Soit l'anneau \( \eZ^2\). L'élément \( (1,0)\) est premier mais pas irréductible.
    \begin{subproof}
        \item[\( (1,0)\) est premier]
            L'élément \( (1,0)\) est non nul; ça c'est pas cher. Pour qu'il soit inversible, il faudrait \( (1,0)(x,y)=(1,1)\). Entre autres, \( 0\times y=1\), ce qui est impossible. Donc il n'est pas inversible.

            Supposons que \( (1,0)\) divise le produit \( (a,b)(c,d)=(ac,bd)\). Alors il existe \( (x,y)\) tel que \( (1,0)(x,y)=(ac,bd)\). Cela signifie que \( x=ac\) et \( 0\times y=bd\). En particulier, soit \( b=0\) soit \( d=0\). Si \( b=0\), nous avons \( (a,b)=(a,0)\) et effectivement, \( (1,0)\) le divise.
        \item[\( (1,0)\) n'est pas irréductible]
            Nous avons \( (1,0)=(1,0)(1,0)\). Donc l'élément \( (1,0)\) est le produit de deux éléments non inversibles.
    \end{subproof}
\end{example}

\begin{example}
    Si \( \eK\) est un corps, l'élément \( XY\) de \( \eK[X,Y]\) n'est pas premier parce que \( XY\divides X^2Y^2\) alors que \( XY\) ne divise ni \( X^2\) ni \( Y^2\).
\end{example}


\begin{proposition}[\cite{ooJHFCooSbHtEC,MonCerveau}, thème \ref{THEMEooVIQIooOcFAQS}]     \label{PROPooZICGooNmblhl}
    Soit un anneau principal \( A\) et un élément \( p\neq 0\) dans \( A\). Nous avons équivalence de :
    \begin{enumerate}
        \item   \label{ITEMooBTEAooWlFUTX}
            \( (p)\) est un idéal premier,
        \item   \label{ITEMooKQRMooBNPDMX}
            \( p\) est un élément premier,
        \item   \label{ITEMooZYYJooCWiBhL}
            \( p\) est un élément irréductible,
        \item   \label{ITEMooHPAIooYoQzqD}
            \( (p)\) est un idéal maximum propre\quext{Ce «propre» n'est pas dans l'énoncé sur Wikipédia. Je ne comprends pas pourquoi, et j'ai posé la question sur la page de discussion.\\\url{https://fr.wikipedia.org/wiki/Discussion:Idéal_premier}}.
    \end{enumerate}
\end{proposition}

\begin{proof}
    En plusieurs implications.
    \begin{subproof}
        \item[\ref{ITEMooBTEAooWlFUTX} implique \ref{ITEMooKQRMooBNPDMX}]
            En plusieurs points.
            \begin{itemize}
                \item La condition \( p\neq 0\) est dans les hypothèses de la proposition.
                \item Si \( p\) était inversible, nous aurions \( (p)=A\) et donc pas que \( (p)\) est un idéal premier.
                \item Soient \( a,b\in A\) tels que \( p\divides ab\). En particulier, \( ab \in (p)\). Mais comme \( (p)\) est un idéal premier, cela implique soit \( a\in (p)\) soit \( b\in (p)\). Donc soit \( p\) divise \( a\) soit \( p\) divise \( b\).
            \end{itemize}
        \item[\ref{ITEMooKQRMooBNPDMX} implique \ref{ITEMooZYYJooCWiBhL}]
            Un anneau principal est intègre; c'est dans la définition \ref{DEFooGWOZooXzUlhK}. Dans un anneau intègre, tout élément premier est irréductible, c'est la proposition \ref{PROPooWMNPooZdvOBt}.
        \item[\ref{ITEMooZYYJooCWiBhL} implique \ref{ITEMooHPAIooYoQzqD}]
            Soit un idéal \( I\) contenant \( (p)\). Vu que \( A\) est principal, \( I\) est engendré par un seul élément; soit \( I=(a)\). Vu que \( p\in I\), l'élément \( a\) divise \( p\). Mais comme \( p\) est un élément premier, \( a\divides p\) implique \( a=p\) ou \( a=1\). Dans le premier cas, \( I=(a)=(p)\), et dans le second cas, \( I=(a)=(1)=A\). Donc \( (p)\) est bien un idéal maximum.

            De plus l'idéal \( (p)\) est propre. En effet avoir \( (p)=A\) dirait en particulier que \( 1\in (p)\), c'est-à-dire qu'il existe \( x\in A\) tel que \( xp=1\). Or \( p\) étant irréductible, il est non inversible.
        \item[\ref{ITEMooHPAIooYoQzqD} implique \ref{ITEMooBTEAooWlFUTX}]
            C'est la proposition \ref{PROPooRUQKooIfbnQX}\ref{ITEMooTFFQooOUajFw}.
    \end{subproof}
\end{proof}

Un exemple d'élément premier non irréductible est \( [4]_6\) dans l'anneau non principal \( \eZ/6\eZ\). Voir \ref{NORMooAXOKooAQMXoB} et le lemme \ref{LEMooZSELooGOFEIz}.

%---------------------------------------------------------------------------------------------------------------------------
\subsection{Anneau noethérien}
%---------------------------------------------------------------------------------------------------------------------------

\begin{definition}      \label{DEFooPWMHooCnrQuJ}
    Un anneau est dit \defe{noethérien}{anneau!noethérien} si toute suite croissante d'idéaux est stationnaire (à partir d'un certain rang).
\end{definition}

Montrer que tout anneau principal est noethérien est le premier pas pour montrer que tout anneau principal est factoriel.

\begin{lemma}       \label{LEMooHQPVooTfkhRV}
    Tout anneau principal\footnote{Définition \ref{DEFooGWOZooXzUlhK}.} est noethérien.
\end{lemma}

\begin{proof}
    Soit \( (J_n)\) une suite croissante d'idéaux et \( J\) la réunion. L'ensemble \( J\) est encore un idéal parce que les \( J_i\) sont emboités. Étant donné que l'idéal est principal nous pouvons prendre \( a\in J\) tel que \( J=(a)\). Il existe \( N\) tel que \( a\in J_N\). Alors pour tout \( n\geq N\) nous avons
    \begin{equation}
        J\subset J_N\subset J_n\subset J.
    \end{equation}
    La première inclusion est le fait que \( J=(a)\) et \( a\in J_N\). La seconde est la croissance des idéaux et la troisième est le fait que \( J\) est une union. Par conséquent pour tout \( n\geq N\) nous avons \( J_N=J_n=J\). La suite est par conséquent stationnaire.
\end{proof}

\begin{example}
    Il y a moyen d'avoir un anneau noetherien non principal. C'est le cas de \( \eZ/6\eZ\) dont nous parlerons dans \ref{LEMooZSELooGOFEIz}.
\end{example}

\begin{theorem}[\cite{FSwlnf}]      \label{THOooANCAooBChmwp}
    Tout anneau principal est factoriel.
\end{theorem}

\begin{example}[\( \eZ\lbrack i\sqrt{ 5 }\rbrack\) n'est ni factoriel ni principal]     \label{EXooYCTDooGXAjGC}
    Vu que \( (i\sqrt{ 5 })^2=-5\), les éléments de \( \eZ[i\sqrt{ 5 }]\) sont les éléments de \( \eC\) de la forme \( a+bi\sqrt{ 5 }\) avec \( a,b\in \eZ\). Nous définissons la \defe{norme}{norme!sur \( \eZ[i\sqrt{ 5 }]\)} sur \( \eZ[i\sqrt{ 5 }]\) par\footnote{C'est le carré de la norme usuelle, mais c'est l'usage dans le milieu.}
    \begin{equation}
        \begin{aligned}
            N\colon \eZ[i\sqrt{ 5 }]&\to \eN \\
            z&\mapsto z\bar z.
        \end{aligned}
    \end{equation}
    Le fait que ce soit à valeurs dans \( \eN\) est un simple calcul :
    \begin{equation}
        N(x+iy\sqrt{ 5 })=(x+iy\sqrt{ 5 })(x-iy\sqrt{ 5 })=x^2+5y^2.
    \end{equation}
    De plus \( N\) est multiplicative : \( N(z_1z_2)=N(z_1)N(z_2)\).

    Nous pouvons maintenant déterminer les inversibles de \( \eZ[i\sqrt{ 5 }]\). Si \( \alpha\) est inversible, alors il existe \( \beta\) tel que \( \alpha\beta=1\). Au niveau de la norme,
    \begin{equation}
        N(\alpha)N(\beta)=1,
    \end{equation}
    ce qui implique que \( N(\alpha)=1\). Or l'équation \( x^2+5y^2=1\) dans \( \eN\) donne \( y=0\), \( x=\pm 1\).

    Au final, les inversibles de \( \eZ[i\sqrt{ 5 }]\) sont \( \pm 1\).

    L'anneau \( \eZ[i\sqrt{ 5 }]\) n'est alors pas factoriel (définition~\ref{DEFooVCATooPJGWKq}) parce que
    \begin{equation}
        2\times 3=(1+i\sqrt{ 5 })(1-i\sqrt{ 5 }).
    \end{equation}
    Cela donne deux décompositions du nombre \( 6\) en produit d'éléments non associés\footnote{Définition~\ref{DefrXUixs}.} (\( 2\) n'est associé qu'à \( 2\) et \( -2\)) parce que les inversibles sont \( 1\) et \( -1\).

    Le fait que \( \eZ[i\sqrt{ 5 }]\) ne soit pas factoriel implique qu'il ne soit pas principal, théorème~\ref{THOooANCAooBChmwp}.
\end{example}

%---------------------------------------------------------------------------------------------------------------------------
\subsection{Anneau euclidien}
%---------------------------------------------------------------------------------------------------------------------------

%///////////////////////////////////////////////////////////////////////////////////////////////////////////////////////////
\subsubsection{Définition et lien avec les autres}
%///////////////////////////////////////////////////////////////////////////////////////////////////////////////////////////

\begin{definition}[\wikipedia{fr}{Anneau_euclidien}{Wikipédia}] \label{DefAXitWRL}
    Soit \( A\) un anneau intègre. Un \defe{stathme euclidien}{stathme euclidien} sur \( A\) est une application \( \alpha\colon A\setminus\{ 0 \}\to \eN\) tel que
    \begin{enumerate}
        \item       \label{ITEMooLVJAooLpjgEz}
            \( \forall a,b\in A\setminus\{ 0 \}\), il existe \( q,r\in A\) tel que
            \begin{equation}
                a=qb+r
            \end{equation}
            et \( \alpha(r)<\alpha(b)\).
        \item
            Pour tout \( a,b\in A\setminus\{ 0 \}\), \( \alpha(b)\leq \alpha(ab)\).
    \end{enumerate}
    Un anneau est \defe{euclidien}{euclidien!anneau} s'il accepte un stathme euclidien.
\end{definition}
Le stathme est la fonction qui donne le «degré» à utiliser dans la division euclidienne. La contrainte est que le degré du reste soit plus petit que le degré du dividende.

\begin{example} \label{ExwqlCwvV}
    Le stathme de \( \eN\) pour la division euclidienne usuelle est \( \alpha(n)=n\). Si \( a,b\in \eN\) nous écrivons
    \begin{equation}
        a=qb+r
    \end{equation}
    où \( q\) est l'entier le plus proche \emph{inférieur} à \( a/b\) (on veut que le reste soit positif) et \( r=a-qb\). Nous avons donc
    \begin{equation}
        r-b=a-b(q+1)<a-b\frac{ a }{ b }=0,
    \end{equation}
    ce qui montre que \( r<b\).
\end{example}

Cet exemple ne fonctionne pas avec \( \eZ\) au lieu de \( \eN\) parce que le stathme doit avoir des valeurs dans \( \eN\). Cela ne veut cependant pas dire qu'il n'existe pas de stathme sur \( \eZ\); cela veut seulement dire que \( \alpha(x)=x\) ne fonctionne pas.

\begin{proposition}[\cite{ooELVSooZIZCRn}]\label{Propkllxnv}
    Un anneau euclidien est principal.
\end{proposition}

\begin{proof}
    Soit \( A\) un anneau principal et \( \alpha\) un stathme sur \( A\). Nous considérons un idéal \( I\) non nul de \( A\). Nous devons montrer que \( I\) est généré par un élément. En l'occurrence nous allons montrer qu'un élément \( a\in I\setminus\{ 0 \}\) qui minimise \( \alpha(a)\) va générer\footnote{Un tel élément existe\dots}. Soit \( x\in I\). Par construction, il existe \( q,r\in A\) tels que \( x=aq+r\) avec \( r=0\) ou \( \alpha(r)<\alpha(a)\). Étant donné que \( x,a\in I\), \( r\in I\). Si \( r\neq 0\), alors \( r\) contredirait la minimalité de \( \alpha(a)\). Donc \( r=0\) et \( x=aq\), ce qui signifie que \( I\) est principal.
\end{proof}

\begin{proposition}     \label{PROPooPJGLooQSrJTU}
    L'anneau \( \eZ\) est principal et euclidien.
\end{proposition}

\begin{proof}
    Nous allons seulement montrer que \( \alpha(x)=| x |\) est un stathme euclidien. Ainsi \( \eZ\) sera euclidien et donc principal par la proposition~\ref{Propkllxnv}.

    D'abord \( \eZ\) est intègre, c'est l'exemple~\ref{EXooLDXRooSxUAXs}.

    La condition \( \alpha(b)\leq \alpha(ab)\) est immédiate : \( | a |\leq | ab |\) pour tout \( a,b\in \eZ\).

    Soient maintenant \( a,b\in \eZ\). Nous définissons \( q_0,r_0\in \eN\) tels que
    \begin{equation}
        | a |=q_0| b |+r_0
    \end{equation}
    avec \( r_0<| b |\). Cela existe parce que \( \alpha(x)=x\) est un stathme sur \( \eN\) par l'exemple~\ref{ExwqlCwvV}.

    \begin{subproof}
        \item[Si \( a>0\) et \( b>0\)]

            Alors \( a=q_0b+r_0\) et le couple \( (q_0,r_0)\) vérifie les conditions de la définition~\ref{DefAXitWRL}\ref{ITEMooLVJAooLpjgEz}.

        \item[Si \( a>0\) et \( b<0\)]

            Alors \( a=-q_0b+r_0\), et le couple \( (-q_0,r_0)\) vérifie les conditions de la définition~\ref{DefAXitWRL}\ref{ITEMooLVJAooLpjgEz}.


        \item[Si \( a<0\) et \( b>0\)]
            Alors \( a=-q_0b-r_0\), et le couple \( (-q_0,-r_0)\) vérifie les conditions de la définition~\ref{DefAXitWRL}\ref{ITEMooLVJAooLpjgEz} parce que
            \begin{equation}
                \alpha(-r_0)=r_0<| b |=\alpha(b).
            \end{equation}
        \item[Si \( a<0\) et \( b<0\)]
            Alors \( a=q_0b-r_0\), et le couple \( (q_0,-r_0)\) vérifie les conditions de la définition~\ref{DefAXitWRL}\ref{ITEMooLVJAooLpjgEz}.

    \end{subproof}
\end{proof}

Nous venons de voir que \( \eZ\) est principal; le lemme suivant nous dit que $\eZ[X]$ n'est pas principal, lui.
\begin{lemma}[\cite{ooRQHSooEBZpKe}]        \label{LEMooDJSUooJWyxCL}
    Si $A$ est un anneau intègre qui n'est pas un corps, alors \( A[X]\) n'est pas principal.
\end{lemma}

\begin{proof}
    Soit un élément non nul \( a\in A\).
    \begin{subproof}
        \newcommand{\foo}{A[X]}
    \item[Un idéal principal contenant \( a\) et \( X\) est {A[X]}]

            Soit \( (P)\) un idéal principal contenant \( a\) et \( X\). Vu que \( a\in(P)\), il existe \( Q\) tel que \( a=QP\). Donc \( P\) divise \( a\) dans \( \eZ[X]\). Les degrés font que \( P\) est un polynôme constant, c'est-à-dire en réalité un élément de \( A\). Soit \( P=k\in A\).

            Vu que \( P\) divise \( X\), nous avons aussi \( X=kQ\) pour un certain \( Q\in \eZ[X]\). Les degrés disent qu'il existe \( k'\in A\) tel que \( Q=k'X\) et donc \( Q=k'X=k'kQ\), ce qui implique que \( kk'=1\). L'idéal engendré par \( k\) contient donc en particulier \( kk'=1\) et donc contient \( A[X]\) en entier :
            \begin{equation}
                1=k'k\in k'(P)=(P).
            \end{equation}

        \item[Si \( (a,X)=\foo\) alors \( a\) est inversible]

            Si \( (a,X)=A[X]\), en particulier, \( 1\in (a,X)\), ce qui signifie qu'il existe des polynômes \( U,V\in A[X]\) tels que
            \begin{equation}
                1=UX+Va.
            \end{equation}
            Nous évaluons cette égalité en \( 0\) : vu que \( (UX)(0)=0\) nous avons \( 1=V(0)a\), ce qui signifie que \( V(0)\) est un inverse de \( A\). Donc \( a\) est inversible.


        \item[Si \( a\) n'est pas inversible alors \( (a,X)\) n'est pas principal]

            Si \( (a,X)\) était principal, alors nous aurions, par ce qui est dit plus haut, \( (a,X)=A[X]\). Mais cette dernière égalité impliquerait que \( a\) est inversible.
    \end{subproof}
    En conclusion, si \( A\) n'est pas un corps, il possède un élément ni nul ni inversible. Dans ce cas, l'idéal \( (a,X)\) n'est pas principal dans \( A[X]\) et nous en déduisons que \( A[X]\) n'est pas un anneau principal.
\end{proof}

\begin{lemma}[Thème~\ref{THEMEooZYKFooQXhiPD}]       \label{LEMooIDSKooQfkeKp}
    Si \( \eK\) est un corps\footnote{Définition~\ref{DefTMNooKXHUd}.}, alors l'anneau \( \eK[X]\) est euclidien et principal.
\end{lemma}

\begin{proof}
    Vu que \( \eK\) est un corps, tous les éléments sont inversibles et le degré donne un stathme par le théorème~\ref{ThodivEuclPsFexf}. L'anneau \( \eK[X]\) est donc euclidien et par conséquent principal (proposition~\ref{Propkllxnv}).
\end{proof}

Dans le théorème~\ref{ThoCCHkoU} nous donnerons une preuve directe du fait que \( \eK[X]\) est principal en montrant que tous ses idéaux sont principaux. Nous y démontrerons donc un peu moins pour un peu plus cher, mais avec le plaisir de ne pas devoir passer par un stathme.

%///////////////////////////////////////////////////////////////////////////////////////////////////////////////////////////
\subsubsection{Un exemple pas trivial}
%///////////////////////////////////////////////////////////////////////////////////////////////////////////////////////////

\begin{example} \label{ExeDufyZI}
    Prouvons que \( \eZ[i\sqrt{2}]\) est un anneau euclidien. Pour cela nous démontrons que
    \begin{equation}    \label{EqOZUIooZGmHWl}
        \begin{aligned}
            N\colon \eZ[i\sqrt{2}]&\to \eN \\
            a+bi\sqrt{2}&\mapsto a^2+2b^2
        \end{aligned}
    \end{equation}
    est un stathme euclidien.

    Soient \( z=a+bi\sqrt{2}\), \( t=a'+b'i\sqrt{2}\). Nous cherchons \( q\) et \( r\) tels que la division euclidienne s'écrive \( z=qt+r\). Soient \( \alpha,\beta\in \eQ\) tels que
    \begin{equation}
        \frac{ z }{ t }=\alpha+\beta i\sqrt{2}.
    \end{equation}
    Nous désignons par \( \alpha+\epsilon_1\) et \( \beta+\epsilon_2\) les entiers les plus proches de \( \alpha\) et \( b\). Nous avons \( | \epsilon_1 |,| \epsilon_2 |\leq \frac{ 1 }{2}\). Nous posons alors naturellement
    \begin{equation}
        q=(\alpha+\epsilon_1)+(\beta+\epsilon_2)i\sqrt{2}
    \end{equation}
    et nous calculons \( r=z-qt\) :
    \begin{equation}
        2b'\epsilon_2-a'\epsilon_1+i\sqrt{2}\big( \epsilon_1b'-a'\epsilon_2 \big).
    \end{equation}
    Nous trouvons
    \begin{equation}
        N(r)=a'^2\epsilon_1^2+4b'^2\epsilon_2^2+2a'^2\epsilon_1^2+2b'^2\epsilon_2^2\leq \frac{ 3 }{ 4 }a'^2+\frac{ 3 }{2}b'^2.
    \end{equation}
    D'autre part \( N(t)=a'^2+2b'^2\), et nous avons donc bien \( N(r)<N(t)\).

    En ce qui concerne la seconde propriété du stathme, un petit calcul montre que
    \begin{equation}
        N(zt)=(a^2+2b^2)(a'^2+2b'^2),
    \end{equation}
    et tant que \( t\neq 0\) nous avons bien \( N(zt)>N(z)\).
\end{example}

Notons en particulier que \( \eZ[i\sqrt{2}]\) est factoriel et principal.

\begin{example}[Décomposition en facteurs irréductibles dans ${\eZ[i\sqrt{2}]}$] \label{ExluqIkE}
    Les éléments inversibles de \( \eZ[i\sqrt{2}]\) sont \( \pm 1\), donc deux éléments \( a\) et \( b\) sont associés (définition~\ref{DefrXUixs}) si et seulement si \( a=\pm b\).

    De plus si \( p\) est irréductible\footnote{Définition \ref{DeirredBDhQfA}}, alors \( -p\) est irréductible. Les éléments irréductibles de \( \eZ[i\sqrt{2}]\) arrivent donc par pairs d'éléments associés. Soit \( \{ p_i \}\) une sélection de un élément irréductible parmi chaque paire. Tout élément \( x\) de \( \eZ[i\sqrt{2}]\) peut alors être écrit \( x=\pm p_1^{\alpha_1}\ldots p_n^{\alpha_n}\). Ce fait va être pratique pour comparer des décompositions en facteurs irréductibles d'éléments.
\end{example}

Le lemme suivant fait en pratique partie de l'exemple~\ref{ExmuQisZU}, mais nous l'isolons pour plus de clarté\footnote{Merci à \href{http://fr.wikipedia.org/wiki/Utilisateur:Marvoir}{Marvoir} pour m'avoir souligné le manque.}.
\begin{lemma}       \label{LemTScCIv}
    Si \( a\) et \( b\) sont deux éléments premiers entre eux de \( \eZ[i\sqrt{2}]\), et s'il existe \( y \in  \eZ[i\sqrt{2}]\) tel que \( ab=y^3\), alors \( a\) et \( b\) sont des cubes (dans \( \eZ[i\sqrt{2}]\)).
\end{lemma}

\begin{proof}
    D'après l'exemple~\ref{ExluqIkE} nous pouvons écrire
    \begin{subequations}
        \begin{align}
            y&=\pm p_1^{\sigma_1}\ldots p_n^{\sigma_n}\\
            a&=\pm p_1^{\alpha_1}\ldots p_n^{\alpha_n}\\
            b&=\pm p_1^{\beta_1}\ldots p_n^{\beta_n}
        \end{align}
    \end{subequations}
    où les \( p_i\) sont les irréductibles de \( \eZ[i\sqrt{2}]\) «modulo \( \pm 1\)» au sens où la liste des irréductibles est \( \{ p_i \}\cup\{ -p_i \}\) (union disjointe). Étant donné que \( a\) et \( b\) sont premiers entre eux, \( \alpha_i\) et \( \beta_i\) ne peuvent pas être non nuls en même temps alors que leur somme doit faire \( 3\sigma_i\). Nous avons donc pour chaque \( i\) soit \( \alpha_i=3\sigma_i\) soit \( \beta=3\sigma_i\) (et bien entendu si \( \sigma_i=0\) alors \( \alpha_i=\beta_i=0\)).

    Étant donné que \( \pm 1\) sont également deux cubes, \( a\) et \( b\) sont bien des cubes.

    Notons que nous avons utilisé de façon capitale le fait que \( \eZ[i\sqrt{2}]\) était factoriel.
\end{proof}


%///////////////////////////////////////////////////////////////////////////////////////////////////////////////////////////
\subsubsection{Équations diophantiennes}
%///////////////////////////////////////////////////////////////////////////////////////////////////////////////////////////
%TODO : il y a une équation diophantienne qui semble pas mal ici : http://fr.wikipedia.org/wiki/Entier_quadratique#x2_.2B_5.y2_.3D_p

\begin{example} \label{ExZPVFooPpdKJc}
    L'équation diophantienne
    \begin{equation}
        x^2=3y^2+8
    \end{equation}
    n'a pas de solutions. En effet si nous prenons l'équation modulo \( 3\) nous obtenons
    \begin{equation}
        [x^2]_3=[3y^2+8]_3=[8]_3=[2]_3.
    \end{equation}
    Or dans \( \eZ/3\eZ\), aucun carré n'est égal à deux : \( 0^2=0\neq 2\), \( 1^2=1\neq 2\) et \( 2^2=4=1\neq 2\).
\end{example}

\begin{example}     \label{ExmuQisZU}
    Résolvons l'équation diophantienne\index{équation!diophantienne}
    \begin{equation}
        x^2+2=y^3.
    \end{equation}
    Une première remarque est que \( x\) doit être impair. En effet si \( x=2k\), nous devons avoir \( y^3\) pair. Mais un cube pair est divisible par \( 8\), donc \( y^3=8l\). L'équation devient \( 4k^2+2=8l^3\), c'est-à-dire \( 2k^2+1=4l^3\). Le membre de gauche est impair tandis que celui de droite est pair. Impossible.

    Nous pouvons écrire l'équation sous la forme \( x^2+2=(x+i\sqrt{2})(x-i\sqrt{2})\). Et nous considérons \( \eZ[i\sqrt{2}]\) muni de son stathme \( N\) donné par \eqref{EqOZUIooZGmHWl}.

    L'élément \( i\sqrt{2}\) est irréductible parce que \( N(i\sqrt{2})=2\), et si nous avions \( i\sqrt{2}=pq\), alors nous aurions \( N(p)N(q)=2\), ce qui n'est possible que si \( N(p)\) ou \( N(q)\) vaut \( 1\).

    Nous prouvons maintenant que les éléments \( x+i\sqrt{2}\) et \( x-i\sqrt{2}\) sont premiers entre eux. Supposons que \( d\) soit un diviseur commun; alors il divise aussi la somme et la différence. Donc \( d\) divise à la fois \( 2x\) et \( 2i\sqrt{2}\).

    Étant donné que \( i\sqrt{2}\) est irréductible et que \( 2i\sqrt{2}=(-i\sqrt{2})^3\), les diviseurs de \( 2i\sqrt{2}\) sont les puissances de \( (-i\sqrt{2})\). Du coup nous devrions avoir \( d=(i\sqrt{2})^{\beta}\) et donc
    \begin{equation}
        x=(i\sqrt{2})^{\beta}q
    \end{equation}
    pour un certain \( q\in\eZ[i\sqrt{2}]\). Dans ce cas nous avons \( N(x)=2^{\beta}N(q)\), mais nous avons déjà précisé que \( x\) ne pouvait pas être pair, donc \( \beta=0\) et nous avons \( d=1\).

    Vu que les nombres \( x\pm i\sqrt{2}\) sont premiers entre eux et que leur produit doit être un cube, ils doivent être séparément des cubes (lemme~\ref{LemTScCIv}). Nous devons donc résoudre séparément \( x\pm i\sqrt{2}=y^3\).

    Cherchons les \( x\) et \( y\) entiers tels que \( x+i\sqrt{2}=y^3\). Si nous posons \( z=a+bi\sqrt{2}\), il suffit de calculer \( z^3\) :
    \begin{verbatim}
----------------------------------------------------------------------
| Sage Version 4.8, Release Date: 2012-01-20                         |
| Type notebook() for the GUI, and license() for information.        |
----------------------------------------------------------------------
sage: var('a,b')
(a, b)
sage: z=a+I*sqrt(2)*b
sage: (z**3).expand()
3*I*sqrt(2)*a^2*b - 2*I*sqrt(2)*b^3 + a^3 - 6*a*b^2
    \end{verbatim}
    En identifiant cela à \( x+i\sqrt{2}\) nous trouvons le système
    \begin{subequations}
        \begin{numcases}{}
            x=a^3-6ab^2\\
            1=3a^2b-2b^3
        \end{numcases}
    \end{subequations}
    où, nous le rappelons, \( x\), \( a\) et \( b\) sont des entiers. La seconde équation montre que \( b\) doit être inversible : \( b(3a^2-2b^2)=1\). Il y a donc les possibilités \( b=\pm 1\). Pour \( b=1\) l'équation devient \( 3a^2-2=1\), c'est-à-dire \( a=\pm 1\). Pour \( b=-1\) l'équation devient \( 3a^2-2=-1\), impossible. En conclusion les possibilités sont
    \begin{subequations}
        \begin{align}
            (x,z)=(-5,1+i\sqrt{2})\\
            (x,z)=(5,-1+i\sqrt{2})\\
        \end{align}
    \end{subequations}
    Le travail avec \( x-i\sqrt{2}\) donne les mêmes résultats.

    Les deux solutions de l'équation \( x^2+2=y^3\) sont alors \( (5,3)\) et \( (-5,3)\).
\end{example}

%///////////////////////////////////////////////////////////////////////////////////////////////////////////////////////////
\subsubsection{Triplets pythagoriciens et équation de Fermat pour \texorpdfstring{$ n=4$}{n=4}}
%///////////////////////////////////////////////////////////////////////////////////////////////////////////////////////////

\begin{definition}
    Les solutions entières (positives) de l'équation \( x^2+y^2=z^2\) sont appelés \defe{triplets pythagoriciens}{triplet!pythagoricien}.
\end{definition}

Ils donnent toutes les possibilités de triangles rectangles dont les côtés ont des longueurs entières.

\begin{definition}
  On dit qu'un triplet pythagoricien est \defe{primitif}{primitif!triplet pythagoricien} si les trois nombres sont premiers dans leur ensemble\footnote{Définition \ref{DefZHRXooNeWIcB}.}.
\end{definition}

Remarquons que cela est équivalent à montrer que les trois nombres sont premiers deux à deux: en effet, si deux parmi \( x\), \( y\) et \( z\) sont divisibles par un nombre, alors tous les trois sont divisibles par ce nombre\footnote{Parce que \( k\) et \( k^2\) ont les mêmes facteurs premiers.}, donc les nombres \( x\), \( y\) et \( z\) sont premiers deux à deux.

\begin{lemma}    \label{LemTripletsPythagoriciensPrimitifs}
  Dans un triplet pythagoricien primitif \( (x, y, z) \), on a toujours $z$ impair et:
  \begin{itemize}
  \item
    soit $x$ impair et $y$ pair;
  \item
    soit $x$ pair et $y$ impair.
  \end{itemize}
\end{lemma}

\begin{proof}
  Remarquons que le fait d'imposer que le triplet soit primitif, interdit aux nombres $x$ et $y$ d'être pairs en même temps. En effet, si c'était le cas, alors \( x^2 \) et \( y^2 \) seraient aussi pairs, donc leur somme \( z^2 \) aussi, d'ou $z$ serait pair et les trois nombres ne seraient pas premiers entre eux.

  Nous montrons à présent que les nombres \( x\) et \( y\) ne sont pas tous les deux impairs. Par l'absurde, si \( x=2a+1\), nous avons \( x^2=4a^2+4a+1\in [1]_4\); de la même manière,  \( y^2 \in [1]_4\). On en déduit alors que \( z^2=x^2+y^2\in [2]_4\). Le nombre \(  z^2\) est donc pair, donc \( z\) est pair : disons \( z=2c\). Alors, \( z^2=4c^2\in [0]_4\). Comme les classes modulo 4 sont disjointes, nous aboutissons à une contradiction.
\end{proof}

\begin{proposition}[Triplets pythagoriciens\cite{fJhCTE,HARRooBvzbXo}]  \label{PropXHMLooRnJKRi}
    Un triplet \( (x,y,z)\in(\eN^*)^3\) est solution de \( x^2+y^2=z^2\) si et seulement s'il existe \( d\in \eN\) et \( u,v\in \eN^*\) premiers entre eux tels que
    \begin{subequations}        \label{subeqLVHFooVgWsFx}
        \begin{numcases}{}
            x=d(u^2-v^2)\\
            y=2duv\\
            z=d(u^2+v^2)
        \end{numcases}
    \end{subequations}
    ou
    \begin{subequations}
        \begin{numcases}{}
            x=2duv\\
            y=d(u^2-v^2)\\
            z=d(u^2+v^2)
        \end{numcases}
    \end{subequations}
    La différence entre les deux est seulement d'inverser les rôles de \( x\) et \( y\).
\end{proposition}

\begin{proof}
    Montrons d'abord que les formules proposées sont bien des solutions; nous vérifions \eqref{subeqLVHFooVgWsFx} :
    \begin{equation}
        x^2+y^2=d^2(u^2-v^2)+4d^2u^2v^2=d^2(u^2+v^2)^2,
    \end{equation}
    qui correspond bien au \( z^2\) proposé.

    Nous allons maintenant prouver la réciproque : toute solution est d'une des deux formes proposées. Déterminer les triplets primitifs suffira parce que si \( (x,y,z)\) n'est pas une solution primitive, alors en posant \( k=\pgcd(x,y,z)\), le triplet \( \big( \frac{ x }{ k },\frac{ y }{ k },\frac{ z }{ k } \big)\) est primitif. Connaissant les triplets primitifs, nous obtenons tous les autres par simple multiplication.

    Soit donc \( (x,y,z)\) un triplet pythagoricien primitif. Sans perte de généralité\footnote{En échangeant les rôles de $x$ et $y$ ici, nous obtenons à la fin la seconde forme des solutions.}, grâce au lemme \ref{LemTripletsPythagoriciensPrimitifs}, nous pouvons supposer  \( x\) est pair tandis que \( y\) et \( z\) sont impairs. Comme \( x^2=(z+y)(z-y)\), nous avons
    \begin{equation}
        \left( \frac{ x }{2} \right)^2=\left( \frac{ z+y }{2} \right)^2\left( \frac{ z-y }{ 2 } \right)^2.
    \end{equation}
    Vu que \( z\) et \( y\) sont premiers entre eux, les nombres \( z-y\) et \( z+y\) sont également premiers entre eux\footnote{Si \( z-y=kn\) et \( z+y=km\), faisant la somme et la différence on voit que \( y\) et \( z\) sont divisibles par \( k\).}. Donc les facteurs premiers (qui arrivent au moins au carré) de \( (x/2)^2\) sont chacun soit dans \( (z+y)/2\) soit dans \( (z-y)/2\). Nous en déduisons que ces derniers sont des carrés d'entiers. Nous posons
    \begin{equation}
        \begin{aligned}[]
            \frac{ z-y }{2}=u^2&&\frac{ z+y }{2}=v^2.
        \end{aligned}
    \end{equation}
    Bien entendu \( u\) et \( v\) sont non nuls parce que nous avons exclu la possibilité de triplets dont un élément serait nul. Avec tout cela nous avons \( (x/2)^2=u^2v^2\) et donc \( x=2uv\) puis par somme et différence :
    \begin{subequations}
        \begin{numcases}{}
            x=2uv\\
            y=v^2-u^2\\
            z=u^2+v^2,
        \end{numcases}
    \end{subequations}
    ce qu'il fallait.
\end{proof}

\begin{remark}
    Les solutions dans lesquelles \( x\), \( y\) ou \( z\) sont nuls sont faciles à classer. La solution \( (1,0,1)\) n'est pas dans les formes proposées. En effet elle ne peut pas être de la première forme : avoir \( y=0\) demanderait qu'un nombre parmi \( d\), \( u\) et \( v\) soit nul, ce qui est interdit. La solution \( (1,0,1) \) ne peut pas non plus être de la seconde forme parce que \( x\) y est pair.
\end{remark}

\begin{proposition}[\cite{fJhCTE}]      \label{propFKKKooFYQcxE}
    Les équations \( x^4+y^4=z^2\) et \( x^2+y^4=z^4\) n'ont pas de solutions dans \( (\eN^*)^3\).
\end{proposition}
\index{équation!diophantienne}

\begin{proof}
  Si la première équation n'a pas de solutions, alors la seconde n'en
  n'a pas non plus parce que \( z^4\) est un carré. Nous nous
  concentrons donc sur l'équation \( x^4+y^4=z^2\) et nous supposons
  qu'il existe au moins une solution dans \( (\eN^*)^3\). Nous en choisissons une \( (x,y,z)\) avec \( z\) minimum (les \( z\) dans différentes solutions étant dans \( \eN\), il en existe forcément un minimum\footnote{Voir quelque chose comme le lemme~\ref{PropQEPoozLqOQ}.}). Du coup, les trois nombres \( x\), \( y\) et \( z\) sont premiers dans leur ensembles parce que une
  division par leur \( \pgcd\) donnerait une nouvelle solution qui
  contredirait la minimalité de \( z\).

    Nous posons \( x^4=\bar x^2\) et \( y^4=\bar y^2\). Ils vérifient
    l'équation \( \bar x^2+\bar y^2=z^2\) et par la proposition
   ~\ref{PropXHMLooRnJKRi}, il existe \( u,v\in \eN^*\) premiers entre
    eux tels que, sans perte de généralité\footnote{En inversant les
      rôles de $x$ et $y$ au besoin.}, on ait
    \begin{subequations}
        \begin{numcases}{}
            \bar x=2uv\\
            \bar y=u^2-v^2\label{eqnFKKKooFYQcxE1}\\
            z=u^2+v^2.\label{eqnFKKKooFYQcxE2}
        \end{numcases}
    \end{subequations}
    Si \( u\) est pair, alors \( v\) est impair (et inversement) parce
    que \( \pgcd(u,v)=1\) Si \( u\) est pair, alors \( u=2a\) et \(
    v=2b+1\), ce qui donne \( \bar y=4a^2-4b^2-4b-1\in[-1]_4\). Or
    nous avons déjà vu qu'un carré est dans \( [0]_4\) ou dans \(
    [1]_4\). Il faut donc que \( u\) soit impair. Le lemme
   ~\ref{LemTripletsPythagoriciensPrimitifs} implique alors que \( v\)
    soit pair.

    De l'équation~\ref{eqnFKKKooFYQcxE1}, nous en déduisons que \(
    v^2+\bar y=u^2\); de plus \( u^2\), \( v^2\) et \( \bar y\) sont
    premiers dans leur ensemble: en effet, $u$ et $v$ sont premiers
    entre eux, et si l'un parmi \( u^2\) et \( v^2\) a un facteur
    commun avec \( \bar y\), alors l'autre doit l'avoir aussi (dans
    une égalité \( a+b=c\), si deux des nombres ont un diviseur
    commun, le troisième l'a aussi). Comme \( \bar y=y^2\), le triplet
    \( (v,y,u)\) est un triplet pythagoricien primitif. Nous
    réappliquons la proposition~\ref{PropXHMLooRnJKRi}, en se
    souvenant que $v$ est pair: il existe donc deux nombres \( r\) et
    \( s\) premiers entre eux tels que
    \begin{subequations} \label{eqnFKKKooFYQcxE3}
        \begin{numcases}{}
            v=2rs\\
            y=r^2-s^2\\
            u=r^2+s^2.
        \end{numcases}
    \end{subequations}
    Avec cela, \( \bar x=2uv=4rs(r^2+s^2)\). Remarquons que les trois
    nombres \( r\), \( s\) et \( r^2+s^2\) sont premiers entre
    eux dans leur ensemble; or, comme \( \bar x\) est un
    carré ces nombres doivent séparément être des carrés :
    \begin{subequations}
        \begin{numcases}{}
            r=\alpha^2\\
            s=\beta^2\\
            r^2+s^2=\gamma^2.
        \end{numcases}
    \end{subequations}
    En mettant les deux premiers dans la troisième, nous avons \( \alpha^4+\beta^4=\gamma^2\). Donc \( (\alpha^2,\beta^2,\gamma)\) est une solution. Nous allons prouver que \( \gamma<z\), ce qui terminera la preuve, puisque $z$ était supposé minimal. Nous avons :
    \begin{align*}
        z&=u^2+v^2&&\text{par~\ref{eqnFKKKooFYQcxE2}}\\
          &=r^2+s^2+4r^2s^2&&\text{par~\ref{eqnFKKKooFYQcxE3}}\\
          &=\gamma^2+4r^2s^2\\
          &> \gamma^2,
    \end{align*}
    et a fortiori \( \gamma<z\).
\end{proof}

%+++++++++++++++++++++++++++++++++++++++++++++++++++++++++++++++++++++++++++++++++++++++++++++++++++++++++++++++++++++++++++ 
\section{L'anneau \texorpdfstring{$ \eZ/6\eZ$}{Z/6Z}}
%+++++++++++++++++++++++++++++++++++++++++++++++++++++++++++++++++++++++++++++++++++++++++++++++++++++++++++++++++++++++++++
\label{SECooSWGKooEeOZTO}

Nous allons donner quelques propriétés de cet anneau, et en particulier voir que dans cet anneau non intègre, la notion d'élément irréductible n'est pas très intéressante.

Voici pour commencer un calcul la table de multiplication de \( A=\eZ/6\eZ\). Pour les multiples de (par exemple) \( [4]_6\) nous écrivons
\begin{equation}
    1\times [4]_6=[4_6]
\end{equation}
et ensuite
\begin{equation}
    2\times [4]_6=[8]_6=[2]_6,
\end{equation}
puis
\begin{equation}
    3\times [4]_6=[2+4]_6=[6]_6=[0]_6,
\end{equation}
et caetera. Le résultat est :
\begin{equation}
\begin{array}{c|c|c|c|c|c|c}
    \times & [0]_6 & [1]_6  & [2]_6  & [3]_6 & [4]_6 & [5]_6  \\
\hline\hline
[0]_6 & 0 & 0 & 0 & 0 & 0 & 0 \\ 
\hline
[1]_6  & 0 & 1 & 2 & 3 & 4 & 5 \\ 
\hline
[2]_6 & 0 & 2 & 4 & 0 & 2 & 4 \\ 
\hline
[3]_6 & 0 & 3 & 0 & 3 & 0 & 3 \\ 
\hline
[4]_6 & 0 & 4 & 2 & 0 & 4 & 2 \\ 
\hline
[5]_6 & 0 & 5 & 4 & 3 & 2 & 1 \\ 
\hline
\end{array}
\end{equation}
Pour ne pas alourdir, nous n'avons pas écrit \( [x]_6\) partout au lieu de \( x\).

\begin{normaltext}[Inversibles]
    Les éléments inversibles de \( \eZ/6\eZ\) sont ceux qui ont un \( [1]_6\) dans leur table de multiplication. Ce sont donc
    \begin{equation}
        U(\eZ/6\eZ)=\big\{ [1]_6,[5]_6 \big\}.
    \end{equation}
\end{normaltext}

\begin{normaltext}[Diviseurs de zéro]
    Les diviseurs de zéro sont ceux qui ont un \( [0]_6\) dans leur table de multiplication, c'est-à-dire
    \begin{equation}
        \big\{ [2]_6,[3]_6,[4]_6 \big\}.
    \end{equation}
\end{normaltext}

\begin{normaltext}[Irréductibles]
    Les irréductibles sont ceux qui ne sont ni inversibles ni produit de deux éléments non inversibles. Les non inversibles sont :
    \begin{equation}
        \big\{ [0]_6,[2]_6,[3]_6,[4]_6 \}.
    \end{equation}
    Ils sont candidats à être irréductibles. Les produits de ces éléments (on oublie les crochets) sont :
    \begin{subequations}
        \begin{align}
            2\times 2&=4\\
            2\times 3&=0\\
            2\times 4&=2\\
            3\times 3&=3\\
            3\times 4&=0\\
            4\times 4&=4.
        \end{align}
    \end{subequations}
    Donc \( [0]_6\), \( [2]_6\), \( [3]_6\) et \( [4]_6\) ne sont plus candidats à être irréductible. Bref, il ne reste aucun candidats.

    L'anneau \( \eZ/6\eZ\) n'a aucun élément irréductible.
\end{normaltext}

\begin{normaltext}[Éléments premiers]       \label{NORMooAXOKooAQMXoB}
    Les éléments non nuls et non inversibles sont \( 2\), \( 3\) et \( 4\).
    \begin{subproof}
    \item[Pour \( 2\)]
        L'élément \( [2]_6\) divise \( 2\), \( 4\) et \( 0\).
        \begin{itemize}
            \item Les \( (a,b)\) tels que \( ab=2\) sont : $(1,2)$, \( (2,4)\) et \( (5,4)\). L'élément \( 2\) divise donc toujours \( a\) ou \( b\).
            \item Les \( (a,b)\) tels que \( ab=4\) sont : $(1,4)$, \( (2,5)\) et \( (4,4)\). L'élément \( 2\) divise donc toujours \( a\) ou \( b\).
            \item Les \( (a,b)\) tels que \( ab=0\) sont : \( (0,x)\),  $(3,2)$ et \( (4,3)\). L'élément \( 2\) divise donc toujours \( a\) ou \( b\). En particulier, \( [2]_6\) divise \( [0]_6\); c'est important.
        \end{itemize}
        Donc \( [2]_6\) est un élément premier.
    \item[Pour \( 3\)]
        L'élément \( [3]_6\) divise \( 3\) et \( 0\).
        \begin{itemize}
            \item Les \( (a,b)\) tels que \( ab=3\) sont : $(1,3)$ et \( (3,5)\). L'élément \( 3\) divise donc toujours \( a\) ou \( b\).
            \item Les \( (a,b)\) tels que \( ab=0\) sont : \( (0,x)\),  $(3,2)$ et \( (4,3)\). L'élément \( 3\) divise donc toujours \( a\) ou \( b\).
        \end{itemize}
        Donc \( [3]_6\) est un élément premier.
        L'élément \( [4]_6\) divise \( 4\), \( 2\) et \( 0\).
        \begin{itemize}
            \item Les \( (a,b)\) tels que \( ab=4\) sont : $(1,4)$, \( (2,5)\) et \( (4,4)\). L'élément \( 4\) divise donc toujours \( a\) ou \( b\).
            \item Les \( (a,b)\) tels que \( ab=2\) sont : $(1,2)$, \( (2,4)\) et \( (5,4)\). L'élément \( 4\) divise donc toujours \( a\) ou \( b\).
            \item Les \( (a,b)\) tels que \( ab=0\) sont : \( (0,x)\),  $(3,2)$ et \( (4,3)\). L'élément \( 4\) divise donc toujours \( a\) ou \( b\).
        \end{itemize}
        Donc \( [4]_6\) est un élément premier.
    \end{subproof}
    Au final, les éléments premiers dans \( \eZ/6\eZ\) sont 
    \begin{equation}
        \big\{ [2]_6, [3]_6, [4]_6  \big\}.
    \end{equation}
\end{normaltext}

Vous noterez que \( \eZ/6\eZ\) a des éléments premiers non irréductibles. Cela est un contre-exemple à la proposition \ref{PROPooZICGooNmblhl} dans le cas d'un anneau non-intègre.


\begin{lemma}[\cite{MonCerveau}]    \label{LEMooZSELooGOFEIz}
    L'anneau \( \eZ/6\eZ\) est noetherien, mais ni intègre ni principal\footnote{Toutes les définitions dans le thème \ref{THEMEooVIQIooOcFAQS}.}.
\end{lemma}

\begin{proof}
    Vu que c'est un anneau fini, toute suite croissante de quoi que ce soit devient stationnaire; donc \( \eZ/6\eZ\) est noetherien.

    Vu que \( \eZ/6\eZ\) a des diviseurs de zéro, il n'est pas intègre. Et vu qu'il n'est pas intègre, il n'est pas factoriel non plus.
\end{proof}
